% STATUT : À RÉDIGER. Matière disponible dans memoire/main.tex (ancien
% document) : le cadre réglementaire est model-agnostique et se réutilise.
\chapter{Le cadre : DORA, Solvabilité II et le risque cyber}
\label{ch:contexte}

\textit{[Chapitre à rédiger. Réutiliser la partie réglementaire de l'ancien document,
qui ne dépend pas du modèle.]}

% PLAN :
%  1. DORA : objet, entrée en application, périmètre, les CINQ piliers
%     (gouvernance, gestion des incidents, tests, tiers, partage d'info).
%     Bien poser la numérotation P1..P5 utilisée dans tout le mémoire.
%  2. Pourquoi les piliers ne sont pas indépendants : ce sont des domaines
%     de CONTRÔLE, une défaillance de l'un dégrade la capacité de l'autre.
%  3. Solvabilité II et le risque opérationnel : où le cyber se loge,
%     ce que la Formule Standard capture (un pourcentage de primes ou de
%     provisions) et ce qu'elle ne capture pas (aucune sensibilité au
%     niveau de maîtrise du risque).
%  4. L'écart que le mémoire vient combler.
